\section{The tensor geometry and the hole relation}\label{sec:geometry}

We construct a triangle-free relation of holes from tensor equations.
The complement of this relation will be our base graph; the moment
arguments in the next sections force its conflict properties.
All vector spaces, linear maps, tensor products and ranks in the construction are over \(\F=\mathbb F_2\).
We use the entrywise pairing of matrices:
\[
 \pair{A}{B}=\sum_{r,c}A_{rc}B_{rc}=\Tr(A\trans B).
\]
Probabilities and entropy take real values.

\subsection{The order of parameters}

We first fix the parameter hierarchy, then define the coefficient
space, its tensor-valued frame embeddings, and the hole equations.
The fixed constants below bound the information retained about a unit
and the number of vectors whose equality will certify a conflict.
Their order ensures that the colour-class scale precedes the graph
threshold, and that the thinning rate precedes the ambient dimension
multiplier.

The construction has one asymptotic parameter \(n\).  Every numerical parameter below, including the dimension multiplier \(M_0\), is fixed before \(n\) tends to infinity.
Set
\begin{equation}\label{eq:early-constants}
 C_0=1000,\qquad g=10^9+1,\qquad
 D=4C_0g,\qquad M=2^{1000},\qquad
 k_{\max}=56(g-1).
\end{equation}
In particular, \(g\) is odd.  The first pin budget and the entropy slack are
\begin{equation}\label{eq:pin-budgets}
 \zeta=\frac{1}{1000k_{\max}},
 \qquad K_1=\left\lceil\frac{4(D+1)}{\zeta}\right\rceil.
\end{equation}
For the phase estimates and a common upper bound on all pin dimensions, define
\begin{equation}\label{eq:law-budgets}
 \begin{gathered}
 r=2K_1+2,\qquad L_0=\lceil100(D+10)\rceil,\qquad
 d_0=4rL_0,\\
 u_0=10(K_1+d_0+1),\qquad
 K=\left\lceil\frac{10(D+K_1+u_0+10)}{\zeta}\right\rceil .
 \end{gathered}
\end{equation}
The budget \(K\) is a conservative upper bound; the construction uses only
the first peeling, whose pin dimension is at most \(K_1\).
The additional probability thresholds in \Cref{sec:phases} use only these early parameters and fixed numerical tolerances.  They do not depend on the selector or channel dimensions.

The early probability margins determine a colour cap \(\eta_{\rm col}>0\)
and then the graph class scale \(a>0\).  Given the graph threshold
\(\beta>0\), choose the endpoint-density bound \(L\) required in
\Cref{sec:graph-transfer}.  Next choose the selector parameters
\(j_*,b\) and the mixer parameters through \(J\) as in
\Cref{sec:moments}.  Choose \(r_0\) sufficiently large after those choices
and all fixed density bounds, and put \(h=1000r_0\).  After \(h\), choose
a sufficiently small thinning rate \(\rho>0\) as in
\Cref{sec:distributions}.  Finally choose a sufficiently large integer
\(M_0\), also relative to \(1/\rho\), and set
\begin{equation}\label{eq:parameter-order}
 \begin{gathered}
 (C_0,g,D,M,k_{\max},\zeta,K_1;\ r,L_0,d_0,u_0,K)\\
 \longrightarrow\ (\eta_{\rm col},a)
 \ \longrightarrow\ (\beta,L)
 \ \longrightarrow\ (j_*,b,\text{mixer parameters},J)\\
 \longrightarrow\ (r_0,h=1000r_0)
 \ \longrightarrow\ \rho
 \ \longrightarrow\ M_0,\qquad
 N=M_0n,\qquad m=2^{C_0gN}.
 \end{gathered}
\end{equation}
We always take \(n\ge 2r_0\) and subsequently increase its lower bound as needed.
The matrices used below have dimensions depending on \(n\); for each \(n\) they are chosen to satisfy \Cref{lem:generic-mixers}, before any law of units is considered.
The complete ledger of inequalities is in \Cref{sec:parameters}.

An \(O(n)\) bit count below has a constant depending only on parameters chosen before \(M_0\).
It can therefore be made a prescribed small fraction of \(N\).
This is different from an \(o(N)\) bound: after \(M_0\) is fixed, \(n/N=1/M_0\) is constant.
Later fixed tests and tolerances may change how large \(n\) must be; they do not change any earlier construction parameter.

\subsection{Moment matrices and cut profiles}

Let
\[
 \cI=\mathbb Z/g\mathbb Z,\qquad
 \cE=\binom{\cI}{2},\qquad
 I_d=\{d+1,\ldots,d+(g-1)/2\}\subset\cI .
\]
Elements of \(\cI\) are \emph{tags}; elements of \(\cE\) are \emph{components}.
The ordinary, or base, variables comprise \(g^2+3\) blocks of \(n\) bits:
\[
 O_{d,t}\quad(d,t\in\cI),\qquad S,\quad \#,\quad Z.
\]
A separate selector has \(b\) bits.  Its values \(s\in\F^b\) are called \emph{labels}, to distinguish them from tags.

The \(O\)-blocks carry testers whose availability depends on the tag,
while \(S\) carries a tester available at every tag.
The additional \(\#\)-block lets the self-Gram form \(E\), defined
below, distinguish different selector labels.
The \(Z\)-block is absent from \(E\); its free coordinates and
change-of-basis symmetry are used in the later parameter tests.

Let
\[
 p_*=\sum_{j=0}^{j_*}\binom bj,\qquad
 \cB=\F^{p_*}\otimes\F^{\,1+(g^2+3)n}.
\]
The first factor is indexed by squarefree selector monomials of degree at most \(j_*\); the second by the constant and the individual base variables.
If \(p_s\) is the vector of selector-monomial evaluations and \(z\) is a base value, write
\begin{equation}\label{eq:point-evaluation}
 v(s,z)=p_s\otimes(1,z)\in\cB.
\end{equation}
Thus the row and column coordinates both represent monomials of base degree at most one and selector degree at most \(j_*\).

Let \(W\subset\cB\otimes\cB\) be the span of the point moments \(v(s,z)v(s,z)\trans\).
For \(l\in\cI\), let \(W_l\) be the span of those point moments for which
\[
 O_{d,t}=0\quad\text{whenever }l\notin I_d.
\]
Let \(W_*\) be defined by requiring every \(O\)-block to be zero.
The cut-profile space is
\begin{equation}\label{eq:cut-profile-space}
 \cX=
 \left\{x=(x_{\{l,l'\}})_{\{l,l'\}\in\cE}:
 x_{\{l,l'\}}=w_l+w_{l'},\quad w_l\in W_l\right\}.
\end{equation}
A tuple \(w=(w_l)_{l\in\cI}\) as in this display is a \emph{representation} of \(x\).

\begin{lemma}\label{lem:cut-kernel}
The intersection of \(W_l\) over any strict majority of tags is \(W_*\).
Two representations of the same cut profile differ by a constant tuple with value in \(W_*\).
\end{lemma}
\begin{proof}
Fix a base coordinate in \(O_{d,t}\).  It is allowed at exactly \((g-1)/2\) tags.
A strict majority therefore contains a tag where that coordinate is disallowed.
Every row and column involving it is zero in \(W_l\) for that tag.
A matrix in the intersection consequently has zero rows and columns at all \(O\)-coordinates.

The coordinate projection that sets all \(O\)-variables to zero sends each point moment to a point moment defining \(W_*\).
Applied on both matrix factors, it fixes the matrix just described.
That matrix is therefore in \(W_*\), proving the nontrivial inclusion.

If \(w,w'\) give the same differences, then
\(w_l+w'_l=w_{l'}+w'_{l'}\) for every pair of tags.
Their difference is a constant tuple.  Its value belongs to every \(W_l\), hence to \(W_*\).
Conversely every such constant tuple lies in the kernel of the cut map.
\end{proof}

In each \(O_{d,t}\)-block and in \(S\), choose the first \(2r_0\) coordinates in disjoint ordered pairs.
The corresponding linear functionals on moments are
\[
 \eta_{d,t}\bigl(v(s,z)v(s,z)\trans\bigr)
   =\sum_{\rho=1}^{r_0}z_{O_{d,t},2\rho-1}z_{O_{d,t},2\rho},
 \qquad
 \eta_S\bigl(v(s,z)v(s,z)\trans\bigr)
   =\sum_{\rho=1}^{r_0}z_{S,2\rho-1}z_{S,2\rho}.
\]
They are well defined by selecting the indicated ordered matrix entries.
Put
\[
 \eta_O=\sum_{d,t}\eta_{d,t},\qquad \eta=\eta_O+\eta_S.
\]
For a cut profile represented by \(w\), define
\begin{equation}\label{eq:cut-functionals}
 \begin{aligned}
 a_t(x)&=\sum_{l,d}\eta_{d,t}(w_l),&
 a(x)&=\sum_t a_t(x)=\sum_l\eta_O(w_l),\\
 b_t(x)&=\sum_{l\ne t}\eta_S(w_l),&
 \chi_*(w)&=\sum_l\eta(w_l).
 \end{aligned}
\end{equation}
The functions \(a_t,a,b_t\) are well defined on \(\cX\).
Indeed, a constant kernel value in \(W_*\) contributes zero to every \(O\)-tester, while it contributes
\((g-1)\eta_S(w_*)=0\) to \(b_t\).
The quantity \(\chi_*(w)\) is used only when a representation has been specified.

Define a symmetric bilinear form \(T^0\) on \(\cX\) by
\begin{equation}\label{eq:T-zero}
 T^0=a\otimes a+\sum_{t\in\cI}(a_t\otimes b_t+b_t\otimes a_t).
\end{equation}
For \(1\le j\le J\) and ordered components \(e,f\), choose endomorphisms \(L_j^{ef},R_j^{ef}\) of \(\cB\).
Their required properties will be proved in \Cref{lem:generic-mixers}.
Set
\begin{equation}\label{eq:T-one}
 \begin{split}
 B^1(x,y)&=\sum_{j=1}^J\sum_{e,f\in\cE}
   \pair{x_e}{L_j^{ef}y_f(R_j^{ef})\trans},\\
 T^1(x,y)&=B^1(x,y)+B^1(y,x),\qquad
 T=T^0+T^1,\qquad
 r(x)=a+T(x,\cdot)\in\cX^*.
 \end{split}
\end{equation}
Here and later the letter \(r\) in \(r(x)\) denotes an affine-gradient map; the early integer \(r\) in \Cref{eq:law-budgets} is its separately specified rank budget.
Since the mixed terms in \(T^0(x,x)\) cancel, \(T^1\) is alternating, and \(a(x)^2=a(x)\), we have
\begin{equation}\label{eq:T-diagonal}
 T(x,x)=a(x)\qquad(x\in\cX).
\end{equation}

The prescribed self-frame Gram matrix is denoted by \(E\).
It is the sum of the ordered tester matrices for \(\eta\) and a matrix \(E^\#\) characterized on evaluation vectors by
\begin{equation}\label{eq:E-sharp}
 v(s,z)\trans E^\#v(s',z')
   =\sum_{\rho=1}^b(s_\rho+s'_\rho)\,z_\#\trans M_\rho z'_\# ,
\end{equation}
where \(M_\rho\) are \(n\)-by-\(n\) matrices chosen in \Cref{lem:generic-mixers}.
The expression is a bilinear form on \(\cB\) because \(j_*\ge1\).
It vanishes when the two point inputs coincide.
Consequently
\begin{equation}\label{eq:E-moments}
 \pair{E}{w}=\eta(w)\qquad(w\in W).
\end{equation}
The \(Z\)-block does not occur in \(E\).
In particular, any invertible linear change of the \(Z\)-variables preserves \(E\); this symmetry will be used in unary probability estimates.

\subsection{Raw vertices}

For each \(e\in\cE\), take copies \(V_e^+,V_e^-\) of \(\F^N\), paired by the usual dot product.
We regard \(V_e^-\) as the dual of \(V_e^+\) through that pairing.
A raw vertex, also called an \emph{orientation}, consists at each component of four maps
\[
 P_e,Q_e:\cB\longrightarrow V_e^+,V_e^-,
 \qquad X_e,Y_e:\F^h\longrightarrow V_e^+,V_e^-,
\]
where the target is the first or second one as appropriate.
They satisfy
\begin{equation}\label{eq:raw-frame-law}
 [P_e,X_e]\text{ and }[Q_e,Y_e]\text{ are injective},\qquad
 P_e\trans Q_e=E,\quad X_e\trans Q_e=0,\quad P_e\trans Y_e=0.
\end{equation}
No value is prescribed for \(X_e\trans Y_e\).

Write \(d=\dim\cB\).
For all large \(n\), \(N\ge2(d+h)\), by the final choice of \(M_0\).
The frame set is then nonempty.
For example, take the plus columns to be the first \(d+h\) coordinate vectors.
Prescribe the full \((d+h)\)-square Gram matrix to be
\(\begin{psmallmatrix}E&0\\0&0\end{psmallmatrix}\).
Lift its columns to minus vectors using the first \(d+h\) coordinates, and append distinct coordinate vectors in a complementary space of dimension \(d+h\).
These minus columns are independent and have the desired Gram values.

Take the uniform distribution on this finite frame set, independently over components.
Let \(\Omega_n\) be the raw vertex space and \(\rawlaw=\rawlaw_n\) this distribution.
We call \(P_e,Q_e\) the primal maps and \(X_e,Y_e\) the channel maps.
The sampling law \(\pi=\pi_n\), defined in \Cref{sec:distributions},
will preserve the \(\rawlaw\)-marginal of all primal maps and restrict
the conditional channel choices.  All identities in this section hold
for every raw vertex, and therefore under either law.
There are at most \(2^{2|\cE|N(d+h)}\) raw vertices, so
\begin{equation}\label{eq:raw-space-size}
 \log_2|\Omega_n|=O(N^2).
\end{equation}

Use the common ambient tensor space
\[
 \cV=\bigoplus_{e\in\cE}V_e^+\otimes V_e^-.
\]
For a raw vertex \(o\), define
\begin{equation}\label{eq:raw-embedding}
 U_o x=(P_{o,e}x_eQ_{o,e}\trans)_{e\in\cE},
 \qquad
 u_o(\xi)=\sum_{e\in\cE}\pair{Y_{o,e}X_{o,e}\trans}{\xi_e}.
\end{equation}
On a simple tensor \(pq\trans\) in component \(e\), the latter pairing is
\[
 u_o(pq\trans)=(Y_{o,e}\trans p)\cdot(X_{o,e}\trans q).
\]
Thus the minus channel is paired with the first tensor factor and the plus channel with the second.
All copies of coefficient spaces attached to different vertices are distinct, although the ambient \(V_e^\pm\) are common.

Each \(U_o\) is injective: a left inverse of \(P_{o,e}\) and the transpose of a left inverse of \(Q_{o,e}\) recover \(x_e\) from \(P_{o,e}x_eQ_{o,e}\trans\).
The two annihilator conditions in \Cref{eq:raw-frame-law} also give
\begin{equation}\label{eq:self-contraction-zero}
 u_oU_o=0.
\end{equation}

\begin{definition}\label{def:hole}
Two raw vertices \(i,j\) have a \emph{hole} between them if there exist \(\lambda_i,\lambda_j\in\cX\) such that
\begin{equation}\label{eq:hole-equations}
 \begin{gathered}
 U_i\lambda_i=U_j\lambda_j,\qquad
 u_jU_i=r(\lambda_i),\qquad u_iU_j=r(\lambda_j),\\
 a(\lambda_i)+a(\lambda_j)=1.
 \end{gathered}
\end{equation}
The middle two identities are identities of linear functionals on the whole cut-profile space.
\end{definition}

\begin{lemma}\label{lem:hole-triangle-free}
The hole relation is symmetric, has no loops, and is triangle-free.
\end{lemma}
\begin{proof}
Symmetry follows by interchanging the endpoints.
A loop would have \(U_i\lambda_i=U_i\lambda_j\), hence \(\lambda_i=\lambda_j\) by injectivity, contradicting the last equation of \Cref{eq:hole-equations}.

Suppose that three distinct raw vertices have holes on all three pairs.
For a directed incidence \(ij\), let \(x_{ij}\) be its witness at \(i\).
For distinct \(i,j,k\), evaluate the gradient identity for \(ij\) on \(x_{ik}\):
\[
 u_j(U_ix_{ik})=a(x_{ik})+T(x_{ij},x_{ik}).
\]
Sum over the six ordered choices of \((i,j,k)\).
For fixed \(j\), the two arguments on the left are equal by the sharing equation for the remaining edge \(ik\), so they cancel.
For fixed \(i\), the two bilinear terms cancel by symmetry of \(T\).
The remaining sum is
\[
 \sum_{\{i,j\}}\bigl(a(x_{ij})+a(x_{ji})\bigr)=1+1+1=1,
\]
a contradiction.
\end{proof}

Sample \(m\) raw vertices independently from the law \(\pi\) defined in
\Cref{sec:distributions}, and form \(G\) on their \emph{positions}: two distinct positions are adjacent when their raw types have no hole.
This gives a finite simple graph even when types repeat.
Equal raw types are adjacent because holes have no loops.
An independent triple of positions would therefore give three distinct raw types forming a hole triangle.
By \Cref{lem:hole-triangle-free},
\begin{equation}\label{eq:sample-independence}
 \alpha(G)\le2,\qquad \chi(G)\ge\lceil m/2\rceil.
\end{equation}
